% TAB-02: generated by harness/analysis/paper_outputs.py; do not edit (every cell's source is in TAB-02.cells.json)
\begin{tabularx}{\linewidth}{l>{\hsize=0.464\hsize\raggedleft\hyphenpenalty=10000\exhyphenpenalty=10000\arraybackslash}X>{\hsize=1.091\hsize\raggedleft\hyphenpenalty=10000\exhyphenpenalty=10000\arraybackslash}X>{\hsize=0.527\hsize\raggedleft\hyphenpenalty=10000\exhyphenpenalty=10000\arraybackslash}X>{\hsize=1.918\hsize\raggedright\hyphenpenalty=10000\exhyphenpenalty=10000\arraybackslash}X}
\toprule
Test & Error-rate ratio & Error-rate ratio 95\% interval & Tie-decided share (\%) & Near-ceiling label \\
\midrule
Test~1 & undefined & comparator error rate at most 4.51 (exact 95\%) & 1.25 & near ceiling (Amendment~1) \\
Test~2 & 11.04 & \mbox{[4.75, unbounded]} & 3.61 & near ceiling (Amendment~1) \\
Test~3 & 2.23 & \mbox{[1.47, 3.92]} & 6.19 & --- \\
Test~4 & 12.33 & \mbox{[5.17, unbounded]} & --- & --- \\
Test~5 & 0.13 & \mbox{[0.00, 0.40]} & --- & --- \\
\bottomrule
\end{tabularx}
% note: Study~1's tests were pre-registered at 24 steps and run at 6 steps (Deviation~1).
% note: Sonnet ran at medium effort, one level below its default, in Tests~3 and 5.
% note: Under the original pre-registered marking rule, which Deviation~2 kept as a sensitivity check, Sonnet-medium's chain score is 55.00\% instead of 80.00\%, and Test~3 would read inconclusive (could not show) instead of worse by more than the margin ($d$~=~+0.30 points, 95\% interval \mbox{[\ensuremath{-}13.58, 14.95]}).
% note: Under the same rule, Test~5's reading is unchanged: shown no worse than Sonnet-medium by more than the margin ($d$~=~+41.25 points, 95\% interval \mbox{[28.75, 52.50]}).
% note: Columns. $k$ is the number of Haiku-medium samples in the vote; raw $k$ is the cost-matched $k$ before the cap (the comparator's mean cost per puzzle over Haiku-medium's mean cost per sample, rounded down); the cap is the most samples a vote may use, 40 for every vote. $p$ (one-sided) is the bootstrap $p$ of the one-sided test of $d$ against the margin (the no-worse-than reading); the Holm threshold is the level that Holm's order sets for that test. The error-rate ratio is the challenger's errors over the comparator's, counted over the puzzles (for a vote, its expected errors); it is pre-specified descriptive, reported beside the test, not part of its reading. The tie-decided share is the mean over puzzles of the share of the vote's draws of $k$ samples in which a tie that included the right answer decided the credit.
% note: Near ceiling (Amendment~1): a label that the pre-registration's Amendment~1 turned on for Tests~1 and 2 because their comparator, Opus-medium, was right on 40 of 40 answers in the length sweep at 6 steps, a pooled whole-answer accuracy of 100.00\%, at or above the pre-set 95\%; it marks a comparator at or near a perfect score. With it, the error-rate ratio (the challenger's errors over Opus-medium's) is reported with its interval next to the no-worse-than reading; where Opus-medium made no errors the ratio is undefined, and its exact 95\% upper error rate, in \%, is given instead.
% note: Format compliance is the share of samples in the required answer format; for Tests~1--3 it is that of the Haiku-medium samples the vote is drawn from. No format-compliance figure is stamped for step voting (Test~4's challenger), so that cell is a dash.
% note: \ensuremath{\leq}~marks a $p$ at the smallest value the bootstrap can return (no resample reached the margin).
